\documentclass[11pt]{article}
\usepackage[margin=0.85in]{geometry}
\usepackage[T1]{fontenc}
\usepackage{amsmath,amssymb,booktabs,longtable,array,xurl,hyperref,upquote}
\hypersetup{colorlinks=true,urlcolor=blue,linkcolor=blue}
\setlength{\emergencystretch}{3em}
\title{Auditing the NRMM Estimator Improvement\newline Against Sharma, Alai and Rajamani (2026)}
\author{Pan Dynamics Collision Avoidance Research}
\date{September 30, 2026}
\begin{document}
\maketitle

\section{Finding}
The current estimator is a useful \emph{structural and maneuvering} improvement,
but the evidence does not support uniformly superior estimation accuracy.
After independent scenario-specific baseline tuning, it wins target-speed
RMSE in three of six scenarios, relative-position RMSE in two of six, and
course RMSE in five of six over the declared post-2-second window. Including
the prescribed initialization transient, it has lower full-run speed RMSE in
all six scenarios. It is substantially better during ego braking and
weaving, while the reference is better on several steady-motion metrics.
Giving both target observers exact ego inputs and the same injection gains
nearly removes the difference. Additional GNSS velocity information and gain
selection must therefore be separated from claims about the target algorithm.


\section{Research question and evidence provenance}
Does the current estimator improve on the NRMM multistage observer in
G.~Sharma, H.~Alai and R.~Rajamani, \emph{Simultaneous ego-vehicle state
estimation and vehicle trajectory tracking using a multistage high gain
observer}, Transportation Research Part C 182 (2026), 105411,
\href{https://doi.org/10.1016/j.trc.2025.105411}{doi:10.1016/j.trc.2025.105411}?
The supplied local PDF was read, including Sections 2--4 and the conclusion;
the university publication record independently confirms its identity.
The experiment concerns the paper's NRMM observer, not its InMM or older RMM
comparators. The paper evaluates a full-scale Chrysler Pacifica with GPS,
IMU and radar. Its original vehicle recordings and executable implementation
were not available here; the data-availability statement offers data on
request. No author was contacted.

Material passport: academic-research-suite, experiment-agent run/validate
workflow; preparation date September 30, 2026; source is an existing local
paper reproduction plus the unchanged production estimator. This is a new
synthetic benchmark and a reproducibility audit, not a reproduction of the
paper's road experiment. The pre-existing September 24 comparison and source
files were present as uncommitted work at intake. This task audited and
extended those sources; it did not newly create every baseline component.

\section{What changed algorithmically}
The target assumption is retained: constant scalar acceleration $A$, constant
sideslip $\beta_C$ (equivalently constant curvature
$\kappa=\sin\beta_C/l_{r,C}$), and strictly positive target speed.
The new coordinates are
\[
 \rho=R_E^T(p_C-p_E),\quad q=R_E^T v_C,\quad s=R_E^T a_C,\qquad
 J=\begin{bmatrix}0&-1\\1&0\end{bmatrix}.
\]
The implemented model satisfies
\begin{align*}
 \dot\rho&=q-v_E-\omega_EJ\rho, &
 \dot q&=s-\omega_EJq, &
 \dot s&=\Phi(q,s)-\omega_EJs,\\
 \Phi&=-\Omega^2q+3A\Omega Jq/\|q\|, &
 A&=q^Ts/\|q\|, &
 \Omega&=(Jq)^Ts/\|q\|^2.
\end{align*}
This is a coordinate reformulation of the retained physical target model,
not a new motion model. Unlike differentiating relative position repeatedly,
it needs neither ego acceleration derivatives nor ego angular acceleration.
It directly addresses the zero-ego-jerk limitation identified in the original
paper's concluding paragraph.

The ego velocity stage now estimates two body-frame components from GNSS
velocity magnitude, accelerometer and gyro readings. It uses a forward-motion
cone and a bounded single-track mismatch
$\omega_E=v_{E,y}/l_E+d_{\rm st}$, rather than deriving inertial velocity
and acceleration from GPS position. A parallel yaw estimate is used only
for inertial outputs. Thus yaw error does not feed back into the eight-state
body-velocity/target core. The full implementation also maintains position,
yaw, predictors, orientation sets and error-bound histories; eight is not
its total software-state count.

\begin{longtable}{@{}p{0.20\linewidth}p{0.35\linewidth}p{0.37\linewidth}@{}}
\toprule
Aspect & Sharma NRMM reproduction & Current estimator \\
\midrule
Target coordinates & $[r_x,\dot r_x,\ddot r_x,r_y,\dot r_y,\ddot r_y]$ & $[\rho;q;s]$ in ego coordinates \\
Ego information & GPS position and gyro for the six-state chain; IMU for yaw/target reconstruction & GNSS velocity, IMU and gyro for body velocity; position is a separate output \\
Ego derivative assumption & $\dot a_x=\dot a_y=\dot\omega_E=0$ in the model derivation & Not needed by the covariant core \\
Additional premises & Acceleration/course heading reconstruction must be informative & Forward cone, calibrated sensors and a quantified single-track mismatch \\
Yaw coupling & Yaw rotates ego velocity into target dynamics & Yaw affects inertial outputs only \\
Target nonlinearity & Positive-speed rational map; reproduction uses a 1 m/s numerical floor & Globally Lipschitz extension equal to the model on its declared domain \\
Certificate & Theorem 1 high-gain condition; eq.~(73) treats yaw rates as fixed parameters & Structured velocity/acceleration sensitivities and a two-stage continuous ISS comparison \\
\bottomrule
\end{longtable}

Removing derivative assumptions while adding sensor and kinematic premises is
not an unconditional weakening of every assumption. Also, Section 8 of
\path{estimator/OBSERVER_ISS_THEORY.md} compares the current certificate with
an earlier yaw-mediated repository design, not directly with Sharma's
Theorem 1. That inequality alone cannot prove that this implementation has
uniformly lower RMSE than the paper.

\section{Baseline fidelity and disclosed implementation choices}
The existing reproduction was checked against eqs.~(26)--(33), (34)--(39),
and (47)--(73). Its normalized LMI minimizes $\lambda_{\max}(P)$ with
$P\succeq I$ and $\lambda=1$; the paper specifies no objective or numerical
gains. ``Certified'' below means $\theta=1.1\max(1,\theta_0)$ using the
printed threshold, not certification of the sampled simulation. ``Matched''
matches the slowest real pole with the current observer; that does not match
all poles or the noise transfer function.

The yaw stage uses a signed angle between estimated inertial acceleration
and measured body acceleration, with course fallback if its estimated
acceleration is too small. Direct visual inspection of PDF page 7 confirms
that printed eq.~(34) lacks the normalization needed for a dimensionless
cosine and uses an unsigned inverse cosine. The signed \texttt{atan2}
implementation repairs these issues; it is not algebraically identical to
the printed expression. Gyro/acceleration thresholds are 0.05 rad/s and
0.5 m/s$^2$, combined by an OR of magnitude tests. These values and the
precise switching convention are reproduction choices, not recovered author
settings. GPS-chain acceleration is a second derivative of position; the
older comparison report's ``third derivative'' description was inaccurate.

Differentiating the companion map under Assumption 1 exposes an omitted
$\omega_EJ\dot v_E$ term in printed eqs.~(62)--(63). The existing finite
difference test independently compares it with an exact physical trajectory.
Both printed and corrected versions are retained in this audit. This is a
formula-level finding; it does not establish what the authors' unavailable
software actually implemented.

Two sampled realizations are retained: zero-order-held positions and an
intersample output predictor. IMU and gyro samples are held over a frame.
RK4 uses at most 5 ms steps and an additional fastest-pole limit for the
Sharma runtime. Acceleration-chain states start at zero. Every method starts
from the same physical position/velocity/target initialization offsets.

\section{Protocol fixed before held-out evaluation}
Six open-loop geometries cover straight oncoming motion, curved crossing,
curved following, braking with a swerve, near-zero relative velocity, and
ego weaving. The target obeys the original Sharma model in every case.
Four cases have constant ego body velocity/yaw rate; swerve and weaving
deliberately violate the original zero-derivative assumption. Durations
are respectively 6, 6, 12, 6.5, 12 and 12 s. No avoidance controller is run.

Sensor draws are identical for all methods in each seed: bounded uniform
noise with vector-norm maxima 0.12 m for GPS position and radar position,
0.02 m/s for GNSS velocity, 0.05 m/s$^2$ for acceleration, and 0.002 rad/s
for gyro. Two-component draws are scaled by $1/\sqrt2$. This evaluates the
repository's sensor contract. The paper instead lists 10 Hz GPS, 100 Hz IMU
and gyro, and 20 Hz radar; neither our synchronized 50 Hz nominal case nor
our synchronized 20 Hz stress case recreates that asynchronous sensor suite.
The IMU specifications in the paper are not interchangeable with the
deterministic noise maxima used here.

\paragraph{Training/test separation.}
The baseline gets 54 alternatives: the prior LMI gain shape or normalized
repeated-pole gains $(3,3,1)$; independent ego and target slow rates
$\{1.5,3,6\}$ s$^{-1}$; yaw gains $\{0.1,1,3\}$ s$^{-1}$.
All retain the paper's model and measurement choices, and use the output
predictor. Empirical alternatives do not claim the nonlinear high-gain
threshold. One \emph{global} candidate is selected using only training
seeds 101--102 on all six scenarios, minimizing the mean of
\[
 \frac{\mathrm{RMSE}_{\rho}}{0.1\,\mathrm m}
 +\frac{\mathrm{RMSE}_{V_C}}{0.5\,\mathrm{m/s}}
 +\frac{\mathrm{RMSE}_{\gamma_C}}{5\,\mathrm{deg}}.
\]
A failed training run incurs infinite loss. No scenario-specific or
metric-specific best baseline is selected after seeing test data. Nominal
test seeds are 201--220. This separates noise realizations, not scenario
geometries: no claim of generalization to unseen roads follows.

A supplemental sensitivity analysis additionally selects one candidate per
scenario using the \emph{same training rows and objective}, writing all six
selections before its test trials. It gives the baseline knowledge of its
scenario and tests whether the global tuning restriction affects the verdict.
This extension was specified during the main campaign after early test
results were visible; the selection computation uses no test errors. Its
results are exploratory robustness evidence, not a preregistered confirmation.
The old matched baseline and the scenario-tuned baseline also receive a
noiseless test, to retain favorable cases for the reference method.

\paragraph{Stress tests and ablation.}
Seeds 201--205 also test 20/100 Hz synchronized sampling, doubled noise,
threefold initialization offsets, a 0.5 s radar gap starting one third into
each scenario, and a constant 0.02 rad/s single-track mismatch (the configured
bound). One noiseless trial per scenario is included. Doubled noise exceeds
the declared sensor bounds and is an empirical stress test only.
The independent target-only ablation gives both algorithms identical exact
ego body velocity, acceleration and yaw rate in the four compatible scenarios.
It includes the old matched gains, training-selected corrected gains and
the current target injection coefficients copied into the companion chain.
This last diagnostic also retains each algorithm's own nonlinear extension;
it is not an assertion that their nonlinear maps or rotating-coordinate
injections are identical. Oracle inputs cannot be deployed as a sensor.

\paragraph{Scoring and uncertainty.}
The primary window is $t\ge2$ s; full-run and final-half errors are also
recorded. Position RMSE uses the Euclidean two-dimensional error; target
speed is the scalar norm, distinguished from vector velocity; course errors
are wrapped before conversion to degrees. Failed/nonfinite runs and position
errors above 1 km are counted explicitly. Timing includes the current
implementation's bound/history work, which the baseline does not implement.
Confidence intervals resample the 20 paired \emph{complete seeds}, not the
correlated time frames, with 10,000 bootstrap draws (RNG 9302026).
Intervals quantify pointwise noise uncertainty within these fixed synthetic
scenarios, not simultaneous coverage or real-road population uncertainty.


\section{Held-out results and the revised verdict}
The main campaign completed 648 training runs, 1,278 full-system test runs,
320 target-only oracle runs and 36 integration-refinement runs in 674.1 s
(excluding MATLAB startup). The supplemental campaign added 120 nominal and
12 noiseless full-system runs. These are 2,414 estimator executions across
all experiment stages, not 2,414 independent trajectories or replicates.

The globally selected candidate is \texttt{repeated\_e6\_t6\_y0.1}:
ego and target gains $(18,108,216)$ in each scalar chain and yaw gain
0.1 s$^{-1}$. Against this one global tuning, the current estimator has
lower mean position, speed and course RMSE in all six nominal scenarios.
Speed reductions are 41--79\%, but the curved-following position reduction
is only 0.000231 m (0.65\%), too small to call a substantial engineering gain.

That global result is \emph{not} the final dominance claim. With per-scenario
training, the reference method achieves lower speed RMSE in three of six
scenarios and lower position RMSE in four of six. The current method retains
lower course RMSE in five of six. Stable-scenario losses and maneuvering
benefits both matter. In the following table, C is the current estimator,
G the globally trained reference, and S the scenario-trained reference.
All entries are means of 20 held-out seed-level RMSE values after 2 s.

\begin{center}\small
\begin{tabular}{@{}lrrrrrr@{}}
\toprule
&\multicolumn{3}{c}{Position RMSE (m)}&\multicolumn{3}{c}{Target speed RMSE (m/s)}\\
Scenario&C&G&S&C&G&S\\
\midrule
Straight oncoming & 0.0267 & 0.0350 & 0.0249 & 0.0666 & 0.1980 & 0.0615 \\
Curved crossing & 0.0267 & 0.0350 & 0.0251 & 0.0660 & 0.2460 & 0.0945 \\
Curved following & 0.0352 & 0.0354 & 0.0250 & 0.1476 & 0.2521 & 0.1350 \\
Swerve and braking & 0.0297 & 0.1087 & 0.0624 & 0.0763 & 0.3558 & 0.5486 \\
Low relative speed & 0.0291 & 0.0354 & 0.0250 & 0.0849 & 0.2004 & 0.0540 \\
Ego weaving & 0.0275 & 0.0448 & 0.0448 & 0.0788 & 0.2745 & 0.2745 \\
\bottomrule
\end{tabular}\end{center}

\begin{center}\small
\begin{tabular}{@{}lrrrl@{}}
\toprule
&\multicolumn{3}{c}{Target course RMSE (deg)}&\\
Scenario&C&G&S&Scenario-trained gain candidate\\
\midrule
Straight oncoming & 0.283 & 3.027 & 0.492 & \texttt{repeated\_e1.5\_t3\_y3} \\
Curved crossing & 0.286 & 4.729 & 3.739 & \texttt{repeated\_e3\_t3\_y1} \\
Curved following & 0.536 & 0.740 & 0.359 & \texttt{repeated\_e3\_t3\_y0.1} \\
Swerve and braking & 0.558 & 9.489 & 7.265 & \texttt{lmi\_e1.5\_t3\_y0.1} \\
Low relative speed & 0.314 & 0.747 & 0.334 & \texttt{repeated\_e1.5\_t3\_y0.1} \\
Ego weaving & 0.357 & 1.437 & 1.437 & \texttt{repeated\_e6\_t6\_y0.1} \\
\bottomrule
\end{tabular}\end{center}

Swerve/braking speed RMSE falls from 3.393 m/s in the old matched reproduction
to 0.356 m/s after global training and 0.549 m/s after scenario training;
the current estimator gives 0.0763 m/s. Scenario training improves the joint
position/speed/course objective, not every component individually.
Thus the defensible speed advantage is about 79--86\% here, not the much
larger ratio against the old gain choice alone. Ego weaving gives
0.0788 versus 0.2745 m/s, a 71\% reduction. Conversely, at low relative speed
the scenario-trained reference gives 0.0540 versus 0.0849 m/s: the current
error is about 57\% larger. No uniform accuracy improvement is supported.

\paragraph{Paired uncertainty.}
The following intervals use the scenario-trained baseline and report
$\Delta=\mathrm{RMSE}_{S}-\mathrm{RMSE}_{C}$ in m/s; positive favors the
current estimator. They are pointwise descriptive intervals conditional on
these six geometries and this gain-selection rule. All trial values and all
three metric intervals for every baseline are retained in CSV.
\begin{center}\small
\begin{tabular}{@{}lrrr@{}}
\toprule
Scenario&Mean $\Delta$&95\% bootstrap interval&Current lower / 20\\
\midrule
Straight oncoming & -0.0050 & [-0.0076, -0.0025] & 4 \\
Curved crossing & 0.0285 & [0.0247, 0.0321] & 20 \\
Curved following & -0.0127 & [-0.0164, -0.0089] & 1 \\
Swerve and braking & 0.4724 & [0.4459, 0.4984] & 20 \\
Low relative speed & -0.0309 & [-0.0326, -0.0292] & 0 \\
Ego weaving & 0.1957 & [0.1914, 0.2008] & 20 \\
\bottomrule
\end{tabular}\end{center}

\paragraph{Initialization versus later tracking.}
Window choice materially changes the ranking. With the prescribed identical
physical initialization offsets, the current method has lower full-run speed
RMSE than the scenario-trained reference in all six cases. The following
planned window-sensitivity check keeps that benefit visible alongside the
post-transient counterexamples. Final-half speed errors preserve the same
three wins and three losses as the primary window. These startup results
are conditional on the chosen initialization, not evidence about every
possible measurement-based initialization procedure.
\begin{center}\small
\begin{tabular}{@{}lrrrr@{}}
\toprule
&\multicolumn{2}{c}{Full-run speed RMSE (m/s)}&\multicolumn{2}{c}{Final-half speed RMSE (m/s)}\\
Scenario&Current&Scenario-trained&Current&Scenario-trained\\
\midrule
Straight oncoming & 0.6149 & 0.8634 & 0.0658 & 0.0581 \\
Curved crossing & 0.4225 & 1.6859 & 0.0656 & 0.0868 \\
Curved following & 0.4272 & 0.6131 & 0.1454 & 0.1188 \\
Swerve and braking & 0.5964 & 2.6785 & 0.0742 & 0.5270 \\
Low relative speed & 0.4595 & 0.5727 & 0.0831 & 0.0502 \\
Ego weaving & 0.3815 & 0.7337 & 0.0815 & 0.2911 \\
\bottomrule
\end{tabular}\end{center}

\section{What the ablation establishes}
With exact ego inputs and the same target injection coefficients, the two
coordinate implementations are almost indistinguishable in these compatible
scenarios. The table reports target-speed RMSE only; position, vector velocity
and course are also retained in \texttt{oracle\_summary.csv}. The oracle
input experiment is diagnostic and has no deployable sensor interpretation.
\begin{center}\small
\begin{tabular}{@{}lrrr@{}}
\toprule
Scenario&Current&Sharma, same injection&Sharma, global training\\
\midrule
Straight oncoming & 0.065586 & 0.065585 & 0.135973 \\
Curved crossing & 0.065603 & 0.065682 & 0.136084 \\
Curved following & 0.147582 & 0.147557 & 0.140679 \\
Low relative speed & 0.084860 & 0.084859 & 0.139920 \\
\bottomrule
\end{tabular}\end{center}

The globally trained reference even has lower oracle speed/course errors on
curved following (0.1407 m/s and 0.502 deg) than the current target observer
(0.1476 m/s and 0.526 deg). This rules out interpreting the old full-system
multipliers as an intrinsic accuracy advantage of the target coordinate
change. The measured gains come from a combination of upstream information,
ego architecture, injection-gain choice and applicability during ego maneuvers.
The experiment does not assign an exact percentage of benefit to each cause.

\section{Stress results, counterexamples and numerical checks}
Against the globally trained baseline, speed and course mean RMSE stay lower
for the current implementation in all 36 noisy stress scenario/variant cells.
This statement does not extend scenario-specific tuning to those stress tests;
that additional campaign evaluated nominal and noiseless sensing only.
Position is not uniformly better: curved following has a slightly smaller
baseline position RMSE at 20 Hz and during radar dropout.
The constant single-track mismatch remains within the declared 0.02 rad/s
bound; no conclusion about large or rapidly changing tire slip follows.

Noiseless matched-gain comparisons retain another counterexample. On curved
following, final-half target course RMSE is approximately 0.08663 deg for the
current estimator and 0.0000180 deg for the matched reference. Holding the
current heading measurement for each 20 ms frame produces a half-frame
rotation lag; $0.15\times0.02/2$ rad is 0.08594 deg. This is a specific sampled
realization limitation, not a flaw in the continuous coordinate identity.
Neither the primary 2 s window nor a final-half window establishes uniform
asymptotic performance for every motion.

The current implementation, globally trained reference and corrected reference
have zero failed full-system runs in 306 each; the scenario-trained reference
has zero in 126. The old matched reference also has zero in 126, and the
threshold-gain hold realization has zero in 120. The threshold-gain output
predictor fails in 77 of 120. These failures concern our sampled realization
and the selected printed threshold, not the paper's continuous-time theorem
or its unpublished experiment settings. Failed runs are counted explicitly;
primary paired intervals compare methods with no failed nominal pairs.

Restoring the omitted companion term does not reverse the measured verdict.
Across global-reference variants, the largest mean metric changes are
0.000561 m in position, 0.01188 m/s in speed, and 0.07897 deg in course.
These are measurable changes and should not be described as exactly zero.

Halving the RK4 step from 5 to 2.5 ms preserves all 36 failure classifications.
For the current method, maximum changes across the six checked seeds are
$3.47\times10^{-9}$ m in position RMSE, $3.00\times10^{-8}$ m/s in speed RMSE,
and $9.98\times10^{-8}$ deg in course RMSE. For the tuned reference and its
corrected variant, the maxima are $3.04\times10^{-6}$ m, 0.00118 m/s and
0.00584 deg. The latter yaw-switch sensitivity is small relative to the
reported full-system differences. This refinement checks numerical sensitivity,
not sampled stability or enclosure of the true physical trajectory.

Nominal mean step times range from 2.024 to 2.284 ms for the current code,
versus 0.291 to 0.319 ms for the globally trained reference. The current code
also computes bounds, orientation sets and histories that are absent from the
reference, so this is implementation cost rather than an isolated comparison
of differential equations. It is approximately seven times slower in these
measurements, and no real-time deadline guarantee is asserted.

\section{Interpretation and research decision}
Retain the covariant formulation and yaw-independent body-velocity core:
they remove a concrete model restriction relevant to collision-avoidance
maneuvers and support explicit disturbance accounting. Describe the work as
an architecture and robustness improvement under a stated sensor/kinematic
contract. The present experiments support maneuvering advantages and reveal
steady-scenario accuracy tradeoffs. They do not support a universal RMSE
improvement, a new target motion model, a real-vehicle performance claim, or
publication novelty beyond the single paper reviewed here.

The next technical priorities are the conservatism/noise tradeoff in the
certificate-selected gains, sample-aligned yaw correction, and validation
with realistic GNSS velocity quality, bias, asynchronous sensing and
single-track mismatch. A sensor-matched hybrid baseline and vehicle data
would be needed to distinguish architecture from additional sensor information
more completely. None of those additional experiments is represented as done.

The prior September 24 report should not supply unqualified improvement
multipliers for a manuscript. Its gain matching does not establish an optimal
reference, its 756 comparisons are not independent experimental replicates,
and its headline does not reflect the favorable per-scenario reference settings
identified here. Historical files were preserved rather than rewritten.

\paragraph{Statistical integrity review (11/11 items).}
Scenario-stratified results avoid an aggregate Simpson reversal; no inference
from scenario means to a road population is made. Selection/Berkson and
collider risks are addressed by reporting all variants and failures, with no
post-outcome trial exclusion. No prevalence/base-rate claim is made. Repeated
measurements are paired with the same initial conditions, so a selected
extreme pre/post improvement is not being claimed. Survivorship is explicit.
All metrics are exported rather than selecting significant endpoints; pointwise
intervals are not sold as multiple-testing-adjusted conclusions. The expanded
scenario-tuning analysis is disclosed as exploratory, addressing forking-path
risk. The controlled algorithm intervention supports comparisons within the
simulator; it does not establish a real-world causal claim or reverse-causal
mechanism. The dominant remaining limitations are model/sensor realism and
baseline-design freedom, not unreported p-values.

\section{Validation, artifacts and reproduction}
Verification status: \textbf{VERIFIED} for deterministic replay and executed
checks; statistical interpretation remains conditional on the protocol above.
The final MATLAB regression run passes \textbf{101/101 tests} across eight
classes: baseline reproduction, fairness interventions, online runtime,
direct velocity, structured high gain, cascade certificate, yaw observer and
gain synthesis. Exact replay of the stored straight-oncoming seed reproduces
the measurements, truth, target states and ego yaw bit-for-bit for all six
methods, excluding wall-clock timing. Independent Python recomputation
confirms all global and per-scenario selections from the training CSV and
confirms that training and test seed sets are disjoint.

MATLAB Code Analyzer checked twelve comparison/test files. It reported only
the host's missing custom-settings-file advisory and used factory defaults;
no source diagnostic was reported. Python compilation and Git whitespace
checks passed. The report was compiled and its tables/readability inspected.
The runtime sources matched the snapshot through experiment completion.
The three comparator algorithms were then moved into the required flat
	exttt{estimator/} layout with path-only updates; their equations and gains
were retained, and the baseline/audit tests and exact replay were rechecked.
Early development failures and the corrected test expectation are retained
in the raw logs; the final test run is wholly successful.

The baseline project revision was
\path{e2813a5f0d6be6ba7f8a25de4c3850bf28e9a7a6} plus the disclosed uncommitted
comparison files. MATLAB is R2026a Update 3. Solver dependencies are the
existing YALMIP revision \path{a30979f7e883c9042766d9e71f4c384d3971fef6}
and SeDuMi revision \path{daef0b7ee87c727a6d97bc1ca892f48b599166d7}; neither
dependency tree is included in the project change. NumPy/Python versions,
source/PDF identities, options and measured completion counts are recorded
in the machine-readable manifests.

Compact evidence is under
\path{report/SHARMA_ESTIMATOR_AUDIT_20260930/}: the protocol, all seed-level
training/test/oracle metrics, both selection tables, summary tables, paired
intervals, numerical checks, test outcomes and source manifest. Full nominal
trajectories remain in the local operational output
\path{simulation_output/sharma_fairness_audit_20260930/traces.mat}.
Standalone plots and the derived report PDF are exported outside the
repository; generated binaries, figures, the literature PDF and unrelated
working-tree changes are not part of the source commit.

From the repository root, with the existing solver dependencies available:
\begin{verbatim}
matlab -singleCompThread -batch "addpath('scripts'); ...
  runSharmaFairnessAudit(OutputDirectory='simulation_output/audit');"
matlab -singleCompThread -batch "addpath('scripts'); ...
  runSharmaScenarioTunedAudit('simulation_output/audit'); ...
  validateSharmaFairnessAudit('simulation_output/audit');"
python3 scripts/analyzeSharmaFairnessAudit.py simulation_output/audit
\end{verbatim}
The ellipses above indicate MATLAB line continuation for readability; the
actual batch commands use the function calls on one line. The plots are
optional via \texttt{--figure-directory} pointing outside the repository.
\end{document}
